% Provenance (never printed): NSI terminale "Algorithmique: methode diviser pour regner (tri fusion, rotation
% d'une image d'un quart de tour)". See sources/book2/BRIEFS.md.
\chapter{Divide and Conquer}\label{ch:g12:divide-and-conquer}

A teacher has $1000$ exam papers to put in alphabetical order. She splits the pile in two and gives each half to a
helper; each helper splits again and passes the halves on, until someone holds a single paper, which is sorted. Then
the sorted piles come back up: at each level, two sorted piles are merged into one by looking only at their top
papers. Nobody ever compares two papers far apart in the order, and the work is done in a small number of rounds.

This chapter turns that idea into a method: split a problem into smaller ones of the same kind, solve them, and
combine their answers. It gives \omterm{def:g12:divide-and-conquer:merge}{merge sort}, which sorts far faster than the sorts of the earlier chapters; a power
computed in a handful of multiplications; and an image turned a quarter in place. It also shows where the method
does not pay.

\begin{recall}
A \omterm{def:g12:recursion:recursion}{recursive function} calls itself on smaller inputs, and is proved by induction (chapter \ref{ch:g12:recursion}).
Binary search halves a sorted list at each step, and its cost is the \omterm{def:g11:cost-of-an-algorithm:halvings}{number of halvings} $h(n)$, the whole number with
$2^h \le n < 2^{h+1}$ (chapter \ref{ch:g11:cost-of-an-algorithm}). Insertion sort makes up to about $n^2/2$
comparisons (the same chapter). A picture is a grid of numbers (Book 1, grade 9).
\end{recall}

\section{Split, solve, combine}\label{sec:g12:divide-and-conquer:method}

\begin{definition}[Divide and conquer]\label{def:g12:divide-and-conquer:dc}
\emph{Divide and conquer}\index{divide and conquer} is a way of designing algorithms in three steps: \emph{divide}
the problem into smaller problems of the same kind; \emph{solve} each of them, by the same method, until they are
small enough to be solved directly; \emph{combine} their solutions into a solution of the whole.
\end{definition}

Binary search is the simplest case: one comparison divides the list into two halves, only one of which can hold the
value; nothing is left to combine. Its cost is the \omterm{def:g11:cost-of-an-algorithm:halvings}{number of halvings}. The algorithms below solve both halves, and
their cost depends on how much the combining step costs.

\begin{notation}[Binary logarithm]\label{not:g12:divide-and-conquer:log2}
For $n > 0$, $\log_2 n$\index{binary logarithm log2} is the number $y$ with $2^y = n$: $\log_2 8 = 3$,
$\log_2 1024 = 10$. It is the logarithm of the mathematics book, in base $2$. For a whole number $n \ge 1$, the \omterm{def:g11:cost-of-an-algorithm:halvings}{number
of halvings} is $h(n) = \lfloor \log_2 n \rfloor$, the whole part of $\log_2 n$.
\end{notation}

\section{Merge sort}\label{sec:g12:divide-and-conquer:merge-sort}

\begin{definition}[Merge, merge sort]\label{def:g12:divide-and-conquer:merge}
To \emph{merge}\index{merge (of sorted lists)} two sorted lists is to make one sorted list of all their elements, by
repeatedly moving the smaller of their two first elements to the output. \emph{Merge sort}\index{merge sort} sorts a
list by splitting it into two halves, sorting each by merge sort, and merging the results; a list of $0$ or $1$
element is already sorted.
\end{definition}

\omcode{code/grade-12/06-divide-and-conquer/listings/merge.py}{1}{18}{Merging two sorted lists, with a count of the
comparisons.}\label{lst:g12:divide-and-conquer:merge}

\omoutput{code/grade-12/06-divide-and-conquer/outputs/merge.txt}{Lengths $4$ and $5$ merged in $7$ comparisons; when
one list runs out, the rest of the other is copied without comparing.}\label{out:g12:divide-and-conquer:merge}

\begin{omfigure}
\begin{tikzpicture}[font=\footnotesize, cell/.style={draw, minimum width=6mm, minimum height=5mm},
    done/.style={cell, fill=black!8, text=black!45}]
  \node[left] at (-0.2,1.4) {\code{a}};
  \foreach \v [count=\i] in {3, 8, 12, 20} \node[\ifnum\i<3 done\else cell\fi] (a\i) at (0.6*\i-0.3,1.4) {$\v$};
  \node[left] at (-0.2,0.4) {\code{b}};
  \foreach \v [count=\i] in {1, 9, 10, 25, 30} \node[\ifnum\i<2 done\else cell\fi] (b\i) at (0.6*\i-0.3,0.4) {$\v$};
  \draw[->, thick, omThm] (a3.north) ++(0,0.4) node[above] {\code{i}} -- (a3.north);
  \draw[->, thick, omThm] (b2.south) ++(0,-0.4) node[below] {\code{j}} -- (b2.south);
  \node[left] at (4.9,0.9) {\code{out}};
  \foreach \v [count=\i] in {1, 3, 8} \node[cell, fill=omDef!10] at (4.6+0.6*\i,0.9) {$\v$};
  \node[cell, fill=omThm!12] (o4) at (6.6,0.9) {$9$};
  \draw[->, omThm] (b2.south east) to[out=-60, in=-90, looseness=0.9] (o4.south);
  \node[font=\scriptsize, align=left, right] at (7.1,0.9) {$12 > 9$:\\ move $9$, then \code{j} moves on};
\end{tikzpicture}

{\small One step of \code{merge}. The grey elements are already in \code{out}, which is sorted and holds the smallest
of \code{a} and of \code{b}. Only \code{a[i]} and \code{b[j]} are compared: the smaller is the next element.}
\end{omfigure}

\omcode{code/grade-12/06-divide-and-conquer/listings/merge_sort.py}{1}{22}{\omterm{def:g12:divide-and-conquer:merge}{Merge sort}: divide, solve both halves,
combine.}\label{lst:g12:divide-and-conquer:merge-sort}

\omoutput{code/grade-12/06-divide-and-conquer/outputs/merge_sort.txt}{Eight numbers in $17$ comparisons; a thousand in
$8705$.}\label{out:g12:divide-and-conquer:merge-sort}

\begin{omfigure}
\begin{tikzpicture}[font=\scriptsize, row/.style={draw, rounded corners=1pt, inner sep=2pt, fill=omDef!8},
    srt/.style={row, fill=omMeth!15}]
  \node[row] (r) at (0,0) {38 27 43 3 9 82 10 5};
  \node[row] (l) at (-3,-0.9) {38 27 43 3}; \node[row] (rr) at (3,-0.9) {9 82 10 5};
  \foreach \t/\x in {{38 27}/-4.5, {43 3}/-1.5, {9 82}/1.5, {10 5}/4.5} \node[row] at (\x,-1.8) {\t};
  \foreach \t/\x in {38/-5.0, 27/-4.0, 43/-2.0, 3/-1.0, 9/1.0, 82/2.0, 10/4.0, 5/5.0} \node[row] at (\x,-2.7) {\t};
  \foreach \t/\x in {{27 38}/-4.5, {3 43}/-1.5, {9 82}/1.5, {5 10}/4.5} \node[srt] at (\x,-3.6) {\t};
  \node[srt] at (-3,-4.5) {3 27 38 43}; \node[srt] at (3,-4.5) {5 9 10 82};
  \node[srt] at (0,-5.4) {3 5 9 10 27 38 43 82};
  \draw[->, black!50] (6.3,-0.2) -- node[right, align=left] {divide} (6.3,-2.6);
  \draw[->, omMeth] (6.3,-2.8) -- node[right, align=left] {merge} (6.3,-5.4);
\end{tikzpicture}

{\small \omterm{def:g12:divide-and-conquer:merge}{Merge sort} on eight numbers: three levels of splitting, then three levels of merging (shaded). Each level of
merging handles all eight numbers once.}
\end{omfigure}

How many comparisons? Take $n = 2^h$. The splitting goes down $h$ levels, to lists of one element. At each level of
merging, the \omterm{def:g12:divide-and-conquer:merge}{merges} together output all $n$ elements, and a \omterm{def:g12:divide-and-conquer:merge}{merge} makes at most one comparison per element output:
at most $n$ per level, so at most $n \cdot h = n \log_2 n$ in all (exercise \ref{exo:g12:divide-and-conquer:8}).
Insertion sort can need about $n^2/2$.

\omcode{code/grade-12/06-divide-and-conquer/listings/count_vs_nlogn.py}{1}{13}{Comparisons counted, against $n h$ and
$n^2/2$.}\label{lst:g12:divide-and-conquer:count}

\omoutput{code/grade-12/06-divide-and-conquer/outputs/count_vs_nlogn.txt}{A sorted input needs exactly $n h / 2$; a
random one a little less than $n h$.}\label{out:g12:divide-and-conquer:count}

\begin{omfigure}
\begin{tikzpicture}
\begin{axis}[omrecord, width=11cm, height=6cm, xmode=log, ymode=log, log basis x=2, xlabel={$n$},
    ylabel={comparisons}, legend pos=north west, legend style={font=\footnotesize}]
  \addplot[omDef, thick, mark=*, mark size=1.4pt] table[x=n, y=merge]{code/grade-12/06-divide-and-conquer/figures/comparisons.dat};
  \addplot[omMeth, dashed, thick] table[x=n, y=nlogn]{code/grade-12/06-divide-and-conquer/figures/comparisons.dat};
  \addplot[omThm, thick] table[x=n, y=half_square]{code/grade-12/06-divide-and-conquer/figures/comparisons.dat};
  \legend{merge sort (random input), $n \log_2 n$, $n^2/2$}
\end{axis}
\end{tikzpicture}

{\small Comparisons of \omterm{def:g12:divide-and-conquer:merge}{merge sort} on random lists of $2$ to $8192$ numbers, from \code{figures/comparisons.py}, both
axes doubling at each mark. The count stays just under $n \log_2 n$; at $8192$ numbers it is $96\,178$, against
$33\,554\,432$ for $n^2/2$.}
\end{omfigure}

Comparisons are not the whole cost: \omterm{def:g12:divide-and-conquer:merge}{merge sort} also copies lists. The clock agrees with the count all the same.

\omcode{code/grade-12/06-divide-and-conquer/listings/measure_sorts.py}{1}{35}{\omterm{def:g12:divide-and-conquer:merge}{Merge sort} and insertion sort,
timed.}\label{lst:g12:divide-and-conquer:measure}

\omoutput{code/grade-12/06-divide-and-conquer/outputs/measure_sorts.txt}{On the author's laptop. Each doubling of $n$
about doubles the time of \omterm{def:g12:divide-and-conquer:merge}{merge sort} and quadruples that of insertion sort.}\label{out:g12:divide-and-conquer:measure}

\begin{omfigure}
\begin{tikzpicture}
\begin{axis}[omrecord, width=11cm, height=5.6cm, xlabel={$n$}, ylabel={milliseconds}, legend pos=north west,
    legend style={font=\footnotesize}, ymin=0]
  \addplot[omDef, thick, mark=*, mark size=1.4pt] table[x=n, y=merge]{code/grade-12/06-divide-and-conquer/figures/measured_sorts.dat};
  \addplot[omThm, thick, mark=square*, mark size=1.4pt] table[x=n, y=insertion]{code/grade-12/06-divide-and-conquer/figures/measured_sorts.dat};
  \legend{merge sort, insertion sort}
\end{axis}
\end{tikzpicture}

{\small Times measured on the author's laptop by \code{figures/measured_sorts.py}, on random lists of $250$ to $8000$
numbers. Insertion sort's curve bends upwards like $n^2$; \omterm{def:g12:divide-and-conquer:merge}{merge sort}'s stays near the axis, at about $8$
milliseconds for $8000$ numbers against about $490$.}
\end{omfigure}

\section{Fast exponentiation}\label{sec:g12:divide-and-conquer:power}

Computing $x^{1000}$ by multiplying by $x$ again and again takes $999$ multiplications. But $x^{1000} =
(x^{500})^2$, and $x^{500} = (x^{250})^2$: one problem of half the size, and one multiplication.

\begin{method}[Fast exponentiation]\label{met:g12:divide-and-conquer:fast-power}
\emph{Fast exponentiation}\index{fast exponentiation} computes $x^n$ from $y = x^{\lfloor n/2 \rfloor}$: $x^n = y^2$
if $n$ is even, $x^n = y^2 \times x$ if $n$ is odd; $x^0 = 1$ and $x^1 = x$. Each step costs one or two
multiplications, and there are $h(n)$ steps: at most $2 \lfloor \log_2 n \rfloor$ multiplications.
\end{method}

\omcode{code/grade-12/06-divide-and-conquer/listings/fast_power.py}{1}{40}{\omterm{met:g12:divide-and-conquer:fast-power}{Fast exponentiation}, recursive, with a
loop, and modulo $m$.}\label{lst:g12:divide-and-conquer:power}

\omoutput{code/grade-12/06-divide-and-conquer/outputs/fast_power.txt}{$3^{1000}$ in $14$ multiplications; the modular
version agrees with Python's own \code{pow}.}\label{out:g12:divide-and-conquer:power}

The loop reads the binary digits of $n$ from the right: \code{square} runs through $x$, $x^2$, $x^4$, $x^8$, \dots,
and \code{result} takes in the powers whose digit is $1$; $13 = 8 + 4 + 1$ gives $x^{13} = x^8 \cdot x^4 \cdot x$.
The modular version keeps every number below $m$ by taking the remainder after each multiplication: $7^{10^{18}}$ has
far too many digits to be written, but its remainder modulo $1\,000\,003$ takes fewer than $120$
multiplications of numbers below $m$. Secure
communication relies on exactly this computation (chapter \ref{ch:g12:secure-communication}).

\section{Rotating an image by quarters}\label{sec:g12:divide-and-conquer:rotate}

Turning a square picture a quarter turn clockwise sends the pixel in row $r$, column $c$ to row $c$, column
$n - 1 - r$. The simplest way writes every pixel into a new grid: it needs a second picture's worth of memory.

\omcode{code/grade-12/06-divide-and-conquer/listings/rotate_copy.py}{1}{20}{A quarter turn into a new
grid.}\label{lst:g12:divide-and-conquer:rotate-copy}

\omoutput{code/grade-12/06-divide-and-conquer/outputs/rotate_copy.txt}{The first row becomes the last
column.}\label{out:g12:divide-and-conquer:rotate-copy}

For a picture of side $2^k$, \omterm{def:g12:divide-and-conquer:dc}{divide and conquer} turns it in place. Cut it into four quadrants. Move each quadrant to
the next place clockwise, top-left to top-right to bottom-right to bottom-left, by swapping four pixels at a time;
then turn each quadrant a quarter, by the same method. Only one pixel is ever held aside.

\begin{omfigure}
\begin{tikzpicture}[font=\footnotesize, x=6mm, y=-6mm]
  \foreach \grid/\dx in {{0,1,2,3,4,5,6,7,8,9,10,11,12,13,14,15}/0, {8,9,0,1,12,13,4,5,10,11,2,3,14,15,6,7}/7,
      {12,8,4,0,13,9,5,1,14,10,6,2,15,11,7,3}/14} {
    \foreach \v [count=\k from 0] in \grid {
      \pgfmathtruncatemacro{\r}{div(\k,4)} \pgfmathtruncatemacro{\c}{mod(\k,4)}
      \pgfmathtruncatemacro{\q}{2*(\v>7)+(mod(\v,4)>1)}
      \pgfmathsetmacro{\s}{{6,18,30,44}[\q]}
      \node[draw, minimum size=6mm, fill=omDef!\s] at (\c+\dx+0.5,\r+0.5) {\v};
    }
    \draw[very thick] (\dx,2) -- (\dx+4,2); \draw[very thick] (\dx+2,0) -- (\dx+2,4);
  }
  \draw[->, thick, omThm] (4.4,2) -- node[above, font=\scriptsize] {move the} node[below, font=\scriptsize]
    {quadrants} (6.6,2);
  \draw[->, thick, omThm] (11.4,2) -- node[above, font=\scriptsize] {turn each} node[below, font=\scriptsize]
    {quadrant} (13.6,2);
\end{tikzpicture}

{\small A $4 \times 4$ grid turned by quadrants. First the four blocks of $2 \times 2$ move round clockwise, each
keeping its inside unchanged (the shades follow each block); then each block is turned on its own. The result is the
output of the copy above.}
\end{omfigure}

\omcode{code/grade-12/06-divide-and-conquer/listings/rotate_quadrants.py}{1}{31}{The quarter turn by quadrants, in
place.}\label{lst:g12:divide-and-conquer:rotate-quadrants}

\omoutput{code/grade-12/06-divide-and-conquer/outputs/rotate_quadrants.txt}{The same grid as the
copy.}\label{out:g12:divide-and-conquer:rotate-quadrants}

The same function turns a real picture. \code{arrow.pgm} is a picture of $16 \times 16$ grey levels in a plain-text
format, a short header and then one number per pixel; the program turns it, writes it to a file, reads it back, and prints both as characters.

\omcode{code/grade-12/06-divide-and-conquer/listings/rotate_image.py}{1}{35}{A picture turned by
quadrants.}\label{lst:g12:divide-and-conquer:rotate-image}

\omoutput{code/grade-12/06-divide-and-conquer/outputs/rotate_image.txt}{The arrow pointing up now points right; the
small mark at the bottom left went to the top left.}\label{out:g12:divide-and-conquer:rotate-image}

\section{When it pays, and when it does not}\label{sec:g12:divide-and-conquer:when}

\omterm{def:g12:divide-and-conquer:dc}{Divide and conquer} pays when the parts are of similar sizes and combining is cheap. If a split leaves one part
almost as large as the whole, as in a split that cuts off one element, the levels are $n$ instead of $\log_2 n$, and
a \omterm{def:g12:divide-and-conquer:merge}{merge sort} built that way is no faster than insertion sort. On very small inputs, the calls and copies cost more
than they save: at $n = 250$ \omterm{def:g12:divide-and-conquer:merge}{merge sort} is only about twice as fast. Real sorting functions, Python's \code{sorted}
among them, \omterm{def:g12:divide-and-conquer:merge}{merge} runs of elements but finish small pieces by insertion.

\begin{remark}
The method also finds the largest sum of consecutive elements of a list: the best block lies in the left half, or in
the right half, or straddles the middle, and the third case is found in one pass from the middle outwards (exercise
\ref{exo:g12:divide-and-conquer:15}).
\end{remark}

\section{Exercises}\label{sec:g12:divide-and-conquer:exercises}

\begin{exercise}[$\star$]\label{exo:g12:divide-and-conquer:1}
\omterm{def:g12:divide-and-conquer:merge}{Merge} $[2, 5, 9]$ and $[1, 6, 7, 11]$ step by step. How many comparisons?
\end{exercise}

\begin{exercise}[$\star$]\label{exo:g12:divide-and-conquer:2}
What are the smallest and the largest numbers of comparisons when merging a list of $3$ elements with one of $4$?
Give an example of each.
\end{exercise}

\begin{exercise}[$\star$]\label{exo:g12:divide-and-conquer:3}
Draw the splits and the \omterm{def:g12:divide-and-conquer:merge}{merges} of \code{merge_sort} on $[5, 2, 4, 6, 1, 3]$.
\end{exercise}

\begin{exercise}[$\star$]\label{exo:g12:divide-and-conquer:4}
List the multiplications \code{power(x, 13)} makes. How many are there?
\end{exercise}

\begin{exercise}[$\star$]\label{exo:g12:divide-and-conquer:5}
Write \code{merge_desc(a, b)}, which \omterm{def:g12:divide-and-conquer:merge}{merges} two lists sorted in decreasing order.
\end{exercise}

\begin{exercise}[$\star\star$]\label{exo:g12:divide-and-conquer:6}
A \code{merge} has no line after its loop. What does it return for \code{[1, 4, 9]} and \code{[2, 3]}? Which tests
would have caught it?
\end{exercise}

\begin{exercise}[$\star\star$]\label{exo:g12:divide-and-conquer:7}
The matrix $\begin{pmatrix} 1 & 1 \\ 1 & 0 \end{pmatrix}^n$ holds the Fibonacci number $\mathrm{fib}(n)$ in its top
right corner. Write \code{fib(n)} with \omterm{met:g12:divide-and-conquer:fast-power}{fast exponentiation} of $2 \times 2$ matrices.
\end{exercise}

\begin{exercise}[$\star\star$]\label{exo:g12:divide-and-conquer:8}
Prove by induction on $h$ that \omterm{def:g12:divide-and-conquer:merge}{merge sort} makes at most $n h$ comparisons on a list of $n = 2^h$ elements.
\end{exercise}

\begin{exercise}[$\star\star$]\label{exo:g12:divide-and-conquer:9}
Turn the grid of letters with rows \code{ab}, \code{cd} a quarter clockwise. Then turn the $4 \times 4$ grid with rows
\code{abcd}, \code{efgh}, \code{ijkl}, \code{mnop} by quadrants, giving the grid after each of the two stages.
\end{exercise}

\begin{exercise}[$\star\star$]\label{exo:g12:divide-and-conquer:10}
In \code{merge_sort}, a student writes \code{mid = len(t) // 2 + 1}. What happens on a list of two elements?
\end{exercise}

\begin{exercise}[$\star\star$]\label{exo:g12:divide-and-conquer:11}
Write a \omterm{def:g12:divide-and-conquer:merge}{merge sort} without recursion: start from runs of one element, and \omterm{def:g12:divide-and-conquer:merge}{merge} the runs two by two until one is
left.
\end{exercise}

\begin{exercise}[$\star\star\star$]\label{exo:g12:divide-and-conquer:12}
Prove that \code{power(x, n)} makes at most $2 \lfloor \log_2 n \rfloor$ multiplications for $n \ge 1$, and exactly
(number of binary digits of $n$ $- 1$) $+$ (number of digits $1$ $- 1$).
\end{exercise}

\begin{exercise}[$\star\star\star$]\label{exo:g12:divide-and-conquer:13}
Write \code{rotate_layers(grid)}, which turns a square grid of any side in place, ring by ring from the outside, four
pixels at a time. Test it against \code{rotate_copy}.
\end{exercise}

\begin{exercise}[$\star\star\star$]\label{exo:g12:divide-and-conquer:14}
Show that for $n = 2^h$, \omterm{def:g12:divide-and-conquer:merge}{merge sort} makes at least $n h / 2$ comparisons on every input: whatever the order, its count
is within a factor $2$ of $n \log_2 n$.
\end{exercise}

\begin{exercise}[$\star\star\star$]\label{exo:g12:divide-and-conquer:15}
Write \code{max_subarray(t)}, the largest sum of a block of consecutive elements of a non-empty list, by \omterm{def:g12:divide-and-conquer:dc}{divide and
conquer}. Test it against the sums of all blocks.
\end{exercise}

\section{Problem: Counting Inversions}\label{sec:g12:divide-and-conquer:problem}

\begin{problem}[{Weekend problem --- how far apart are two rankings?}]\label{pb:g12:divide-and-conquer:1}
Two judges rank the same items. Writing the second ranking in terms of the first's positions gives a list; each pair
of items the judges order differently is an \emph{inversion} of that list, a pair of positions $i < j$ with
$t[i] > t[j]$. The program counts inversions by trying every pair, and by a \omterm{def:g12:divide-and-conquer:merge}{merge sort} that counts them on the way.

\omcode{code/grade-12/06-divide-and-conquer/listings/inversions.py}{1}{37}{Inversions, by brute force and by
merging.}\label{lst:g12:divide-and-conquer:inversions}

\omcode{code/grade-12/06-divide-and-conquer/listings/inversions.py}{39}{48}{Two judges, then two random
rankings.}\label{lst:g12:divide-and-conquer:inversions-run}

\omoutput{code/grade-12/06-divide-and-conquer/outputs/inversions.txt}{Two judges of five, then two random rankings of
$100\,000$ items.}\label{out:g12:divide-and-conquer:inversions}

\omcode{code/grade-12/06-divide-and-conquer/listings/measure_inversions.py}{1}{20}{The two methods,
timed.}\label{lst:g12:divide-and-conquer:measure-inversions}

\omoutput{code/grade-12/06-divide-and-conquer/outputs/measure_inversions.txt}{On the author's laptop. Brute force is
not even tried at $100\,000$.}\label{out:g12:divide-and-conquer:measure-inversions}

\medskip\noindent\textbf{Part I --- Rankings.}
\begin{enumerate}
  \item In the example of five, which pairs do the judges order differently?
  \item What list does \code{compare} give for two identical rankings, and how many inversions? For two opposite
    rankings of $n$ items?
  \item Why is the number of inversions at most $n(n - 1)/2$?
  \item Two random rankings: what share of the pairs would you guess to be inverted? Compare with the output.
\end{enumerate}

\medskip\noindent\textbf{Part II --- Brute force.}
\begin{enumerate}[resume]
  \item How many pairs does \code{inversions_brute} examine for $n$ items?
  \item From the measured times, by how much does the brute-force time grow when $n$ doubles? Why?
  \item Estimate its time for $n = 100\,000$ from the time at $8000$.
\end{enumerate}

\medskip\noindent\textbf{Part III --- Counting while merging.}
\begin{enumerate}[resume]
  \item Split the inversions of a list into three kinds, by where their two positions lie.
  \item Which two kinds do the recursive calls count?
  \item Why can the third kind be counted on the two sorted halves, although sorting changed the positions?
  \item Explain the loop over \code{right}: why is \code{len(left) - i} the number of inversions with \code{b}?
  \item Why does \code{i} never go back, and what does that make the cost of the loop?
  \item Prove, by induction on the length, that \code{sort_count} returns the number of inversions.
\end{enumerate}

\medskip\noindent\textbf{Part IV --- Cost and speed.}
\begin{enumerate}[resume]
  \item How many comparisons does \code{sort_count} make, roughly, for $n = 2^h$?
  \item Read the measured ratio of the two methods at each $n$. How does it grow?
  \item From your estimate of question 7, what would the ratio be at $100\,000$?
  \item From $n = 1000$ to $n = 8000$, compare the growth of the measured ratio with that of $n / \log_2 n$.
  \item The brute-force method is still useful. For what?
  \item The tests check the count at $100\,000$ against no brute force. What do they check instead?
  \item State the result: the inversions between the two rankings of $100\,000$, and the measured ratio at $8000$.
\end{enumerate}
\end{problem}
